% Audit proof supplement; not an assembled Paper III.
% GPT-6 Astra (OpenAI, ChatGPT Pro mode), effort Pro, 17 November 2026.
\begin{lemma}[The actual order-four torsion transgression]
\label{mapsaudit:order-four-torsion-transgression}
Use the order-four filling and the marked matrices of MAPS Theorem F.2.6.
Let $T=\bigwedge^{q-1}D^t$ on the integral vertical cochains of the
five-torus. If $w$ represents a torsion element of
$\operatorname{coker}(T-I)$, its point-supported degree-$q$ class has
Leray differential equal, up to the one boundary-orientation sign, to
the image of $w$ in the global degree-$(q-1)$ coinvariants. There is
no multiplicity-four factor. For each of $q=2,3,4$, that image is
the nonzero class in the global $\mathbb Z/2$ quotient.
\end{lemma}
\begin{proof}
In the original integral marking the required matrices are
\[
D=\begin{pmatrix}
0&0&0&1&0\\-1&-1&-1&-1&-1\\0&0&0&0&1\\
1&0&0&0&0\\0&1&0&0&0
\end{pmatrix},\qquad
P=\begin{pmatrix}
-1&0&-1&0&1\\-1&2&-1&1&3\\-1&1&-1&2&3\\
0&-1&0&-1&-1\\1&-1&2&-1&-3
\end{pmatrix}.
\]
Direct multiplication gives $\det P=\pm1$ and $DP=PC$, where
$C$ cyclically permutes the first four basis vectors and negates the
fifth. Thus $P$ changes $D$ to a four-cycle permutation of four circle
coordinates and inversion of the fifth.
Equip each circle with its one-vertex, one-edge cellular structure
and use the product structure. The resulting automorphism is cellular,
acts on degree-$j$ integral cochains by $\bigwedge^jD^t$, and its
fourth power is the identity as a map, not only up to homotopy.
The product cellular differential of the five-torus is zero.
Consequently the mapping-torus cellular cochain complex is the cone
of $D^*-I$: a suspension class in degree $q$ is represented by
the cochain $w$ on its base edge times the vertical
degree-$(q-1)$ cells, modulo $(T-I)$ of the vertical cochains.
This is an actual cellular representative of the local torsion
class; no splitting of a Wang exact sequence is being assumed.

The finite normal line is topologically trivial in the supplied
marked quotient, as follows directly from its affine action. Write
the invariant torus coordinate as $x_0\in\mathbb R/\mathbb Z$ and
use the generator
\[
 (x_0,y,z)\longmapsto(x_0-1/4,Dy,iz).
\]
The function $z_{\mathrm{tr}}=e^{2\pi i x_0}z$ is invariant, and so gives a
nowhere-zero smooth normal trivialization. On the boundary put
$a=-4x_0\bmod1$ and $k=\arg(z_{\mathrm{tr}})/(2\pi)\bmod1$.
The boundary is therefore $B_4\times S^1_k$, and the argument
of the original base coordinate $z^4$ is
\[
             \theta=a+4k.
\]
On the base two-torus choose the new integral generators
\[
 e_\theta=(1,0),\qquad e_h=(-4,1)
\]
in the old $(a,k)$ marking. Their determinant is one. The vertical
monodromy along $e_\theta$ is $T$, and along $e_h$ it is
$T^{-4}=I$. Thus the fibre over $\theta$ is the original
six-torus, with its five vertical directions and the primitive
circle $h$. In particular $e_\theta$ goes around the suspension
circle once; confusing it with a fourfold traverse would introduce
an incorrect factor.

The pullback of a cellular one-cochain with values in a local
system is computed by following its oriented edges and transporting
each value to the starting fibre. Applied to the word
$h=a^{-4}k$, this gives
\[
 (w,0)_{(a,k)}\longmapsto
 \left(w,-(T^{-1}+T^{-2}+T^{-3}+T^{-4})w\right)_{(\theta,h)}.
\tag{FT.1}
\]
One can also obtain this formula by writing the four successive
fundamental prisms: reversal of an $a$ edge contributes $-T^{-1}$,
and the next three reversals contribute $-T^{-2},-T^{-3},-T^{-4}$.
The final $k$ edge contributes zero. Since $T^4=I$, the sum in
\textup{(FT.1)} is $N_T=I+T+T^2+T^3$.

A torsion coinvariant has $w=(T-I)v$ for some rational vertical
cochain $v$. Hence
\[
 N_Tw=N_T(T-I)v=(T^4-I)v=0
\]
as an integral cochain. Formula \textup{(FT.1)} therefore says
that the restriction to the six-torus fibre is zero, while the
first boundary Wang component is exactly $[w]$. This calculation
is unaffected by the ambiguity $w\mapsto w+(T-I)b$: its change is
a coboundary on the suspension circle. The cellular construction
preserves vertical degree, so there is no omitted component in a
lower vertical degree. If the inverse monodromy convention is used,
the same argument replaces $w$ by $-T^{-1}w$, which is $-w$ in
coinvariants and hence the same element in the order-two target.

Use the two-open-cover \v{C}ech resolution for the base, with one
open set containing the finite centre and the other its complement,
and the cellular fibre resolution just specified. The local class
has zero fibre restriction on their annular intersection. Its next
component is the boundary one-cocycle $w$ computed above. The
connecting homomorphism identifies this with its class in degree-two
base cohomology. For the punctured sphere, that group is the global
coinvariant quotient: transporting a coefficient along the two
ordinary generators imposes exactly $(\bigwedge^{q-1}A-I)$ and
$(\bigwedge^{q-1}B-I)$. Thus the two-step differential of this
actual \v{C}ech--fibre complex is precisely
\[
 d_2(w)=\pm[w]\in
 \frac{\bigwedge^{q-1}\mathbb Z^6}
 {(\bigwedge^{q-1}A-I)\bigwedge^{q-1}\mathbb Z^6+
  (\bigwedge^{q-1}B-I)\bigwedge^{q-1}\mathbb Z^6}.
\]
The finite modifications of stalks do not alter this degree-two
base cohomology: the kernel and cokernel sheaves of the comparison
with the invariant-extension sheaf are supported at finitely many
points and have no positive cohomology.

For completeness the supplied local representatives, in increasing
lexicographic exterior coordinates of the vertical five-lattice, are
\[
\begin{array}{c|l}
q=2&(0,0,0,-1,-2),\\
q=3&(0,0,0,0,0,0,0,1,0,-1),\\
q=4&(0,0,0,0,0,0,1,0,0,0).
\end{array}
\]
For each row the norm is zero; twice the vector lies in the integral
image of $T-I$, and the vector itself does not. Its embedding in
the six-lattice fills the coordinate slots not containing the first
basis vector and is zero on the other slots. The global relation
matrix has Smith factors all one except for a single two; adjoining
that embedded vector changes every Smith factor to one. It follows
that its residue generates the global $\mathbb Z/2$. These are
independent integral lattice tests in the accompanying checker, in
addition to the original exact Smith computation. The geometric
identification of the differential was proved above and is not an
inference from the arithmetic tests. The three torsion differentials
of F.2.6 are therefore isomorphisms.
\end{proof}

\paragraph{Dependency scope.}
This proof verifies the formerly unexamined finite-boundary step of
F.2.6. It uses the actual affine filling, its integral marking, the
primitive maximal cusp specialization, and the other already stated
Leray groups of that theorem. It does not replace the separate
existence/properness proof for the rank-three Kuhn cusp or the later
holomorphic period construction.
